\chapter{Fake News}
\label{ch:02_fake-news}

\begin{universityTask}
    Die Nachrichten, die mithilfe von Social Bots verbreitet werden, sind in der Regel sprachliche Nachrichten. Einen besonderen Stellenwert nehmen dabei die oft als \q{Fake News} bezeichneten Lügen und Gerüchte ein.

    \vspace{0.5cm}

    \begin{enumerate}[label=(\alph*)]
        \item Wodurch unterscheidet sich sprachliche von nichtsprachlicher Kommunikation?
            \universityPoints{4}
        \item Was unterscheidet Kommunikation von nichtkommunikativen Interaktionsformen?
            \universityPoints{4}
        \item Unter welchen Bedingungen handelt es sich bei einem Sprechakt um eine Lüge?
            \universityPoints{2}
        \item Unter welchen Bedingungen ist ein Sprechakt wahr?
            \universityPoints{2}
        \item Was versteht man überhaupt unter einem Sprechakt? Und wie hängt der Begriff mit dem des Sprachspiels zusammen?
            \universityPoints{4}
        \item Durch welche (an sich unvollständigen) Teilhandlungen wird der Sachbezug eines Sprechaktes durchgeführt? Was ist jeweils deren Funktion?
            \universityPoints{3}
        \item Geben Sie Interaktions-, Selbst- und ggf. Sachbezüge der folgenden drei Äußerungen an (denken Sie dabei vor allem auch an die korrekte Form, in der die drei Bezüge stets angegeben werden sollten):
        \begin{enumerate}[label=\roman*.]
            \item \qit{Fünfhundert Frauen wurden gestern in Köln von Asylanten vergewaltigt.}
            \item \qit{Vermisst wird Marie, 6 Jahre, zuletzt gesehen mit zwei Flüchtlingskindern.}
            \item \qit{Ich verspreche dir, du machst garantiert 200 \% Gewinn pro Jahr mit den Aktien.}
        \end{enumerate}
            \universityPoints{12}
        \item Geben Sie die Teilhandlungen der Proposition an, die den Sachbezug von Äußerung i) in Teilaufgabe g) bildet?
            \universityPoints{3}
        \pagebreak
        \item Skizzieren Sie die Pragmatik eines Gerüchtes: Welche Handlungen, die mit dem Sprechakt, mit dem das Gerücht erzeugt wird, in Verbindung stehen, sind bei Gerüchten besonders relevant?
            \universityPoints{4}
    \end{enumerate}
\end{universityTask}

Dieser Text enthält die Lösung der obigen Aufgabenstellung.